\input{preamble}

\title{Section 3.2 --- Root Properties}

\begin{document}
\maketitle
\section{Solving Equations with Root Properties}
Solve each equation using the appropriate root property. In this section, all roots will be integers. If there are multiple solutions, list them all. If there is no solution, write ``no solution.''
\begin{smenum}
\mitemxxxx
{$x^3=64$}
{$x^2=25$}
{$x^4=81$}
{$x^7=-128$}
\mitemxxxx
{$4x^2+7=43$}
{$-2x^3-21=33$}
{$2x^2+25=7$}
{$5x^3=625$}
\mitemxxxx
{$(2x-3)^2=25$}
{$(3x-1)^3=0$}
{$(4x+5)^4=1$}
{$(4-x)^5=243$}
\mitemxx
{$11(x+5)^2-4=95$}
{$-4(3-x)^4+17=-47$}
\mitemxx
{$-4(3-x)^3-81=175$}
{$2(3x-5)^5-17=-15$}
\end{smenum}
\section{Simplifying Roots}
Simplify each root.
\begin{smenum}
\mitemxxxx
{$\sqrt{50}$}
{$\sqrt[3]{56}$}
{$\sqrt{216}$}
{$\sqrt[3]{336}$}
\end{smenum}
Simplify each fraction. If necessary, rationalize the denominator. Also simplify the radical in the numerator, if possible.
\begin{smenum}
\mitemxxxx
{$\sqrt{\dfrac{5}{16}}$}
{$\sqrt[4]{\dfrac{11}{81}}$}
{$\sqrt[3]{\dfrac{25}{27}}$}
{$\sqrt{\dfrac{18}{49}}$}
\mitemxxxx
{$\sqrt[4]{\dfrac{7}{27}}$}
{$\sqrt[3]{\dfrac{49}{72}}$}
{$\sqrt{\dfrac{5}{200}}$}
{$\sqrt[3]{\dfrac{125}{12}}$}
\end{smenum}
\section{Simplifying Solutions of Equations}
Solve each equation. Simplify your answer as much as possible. In this section, your answers may or may not be rational. If there are two solutions, list them both using $\pm$ notation and as two separate solutions.
\begin{smenum}
\mitemxxx
{$(x-5)^2=7$}
{$2(x-3)^3+9=5$}
{$5(2x-8)^2+7=47$}
\mitemxxx
{$\frac{1}{2}(x-7)^3-5=7$}
{$3(x-5)^2=20$}
{$4(2x+1)^3=27$}
\end{smenum}
\end{document}
